\section{Introduction}

Unmanned aerial vehicles (UAVs) expands its applications in a diversified area such as surveying and mapping \cite{ref01}, search and rescue \cite{ref02} and many other applications \cite{ref03}. With the recent advancements of UAV technology instead of using single-UAVs, usage of multi-UAVs has been evolving and identifies as a highly interested field in research. The high interest in multi-UAV systems have several advantages comparing to single-UAV systems. Time efficiency is increased by reducing the operations time of the mission, by allocating distributed tasks to each UAV to complete the mission. This is immensely helpful in activities such as target search and urgent detection of nuclear radiation in a disaster \cite{ref03}. Use of multi-UAVs becomes more cost effective in the long run than using a single UAV. This reduces the power consumption costs and becomes cost effective in applications related to heavy weight. Multi-UAV systems helps in fulfilling simultaneous actions in different geographical locations at the exact same time where a single UAV is incapable. Multi-UAV systems act as a solution for fault tolerance applications. Loss of a single UAV in a multi-UAV system doesn't explicitly state the end of a mission. The system compensate the mission using the remaining UAVs. With these advantages, several limitations could be identified in research regarding to multi-UAV systems. Legal restrictions are imposed in some jurisdictions which does not permit to fly in swarms \cite{ref03}.Also the safety concerns regarding safe piloting among swarms causes a severe problem. Such limitations paves path to make research strong on simulation frameworks regarding multi-UAV systems.

Identifying this major problem, simulation tools have been built to mitigate the overheads of real world testing, many such existing simulation tools fail at instances because they often focus solely on either the accurate modelling of UAV operations \cite{drons3} or communication and network dynamics \cite{gazebons3}. The built framework addresses the challenge of accurately modelling complex, coupled systems which include physical, network and easily controllable systems. Specifically, there is a lack of readily available, accessible, modular and extensible open-source frameworks capable realistically simulating the dynamics of multi-UAV systems performing distributed tasks, especially those involved onboard data processing and clustering. Therefore, there exists a clear need for a comprehensive multi-UAV system.

Hence, a multi-UAV simulation framework which is expandable and capable of handling multi-UAV research has been built integrating ArduPilot, an open source UAV flight controller, Gazebo, a 3D environment simulator to simulate the physical environment, ns3 a network simulator to simulate the communication network and ROS2 as the middleware.

The remainder of this paper is organized as follows. Section II reviews related work, Section III details the framework architecture, Section IV describes how the network model is validated, Section V compares edge vs. ground processing, Section VI describes how the HITL validation is taken place, Section VII discusses about the limitations and future work and the last section concludes the paper.
